\chapter{Theoretische und technische Grundlagen}

\section{Eingebettete Systeme \& STM32N6570DK}
\subsection{Begriffsbildung und Merkmale}
\subsection{Mikrocontroller Architektur}
\subsection{Peripherieschnittstellen}
\subsection{Speicherarchitektur}
\subsection{Clock-Systeme}
\subsection{Bare-Metal-Programmierung}
\subsection{Hardware-Beschleunigung}

\section{Deutsches Einhand-Fingeralphabet}
\subsection{Gebärdensprache als eigenständige Sprache}
\subsection{Das Fingeralphabet}
\subsection{Das deutsche Einhand-Fingeralphabet}
\subsection{Handpositionen und Fingerkonfigurationen}

\section{Grundlagen des maschinellen Lernens}
\subsection{Begriffsbildung}
\subsection{Supervised Learning}
\subsection{Klassifikation}
\subsubsection{Sparse Categorical Crossentropy}
\subsection{Metriken}
\subsection{Overfitting und Gegenmaßnahmen}
\subsection{Parameter vs Hyperparameter}
Text.

\section{Neuronale Netze zur Klassifikation}
\subsection{Künstliche Neuronale Netze}
\subsection{Aktivierungsfunktionen}
\subsection{Architektur: Multi-Layer Perceptron (MLP)}
\subsection{Architektur: Convolutional Neural Networks (CNN)}
\subsection{Backpropagation und Gradient Descent}

\section{Handdetektion und Handlandmarks}
\subsection{Zweistufige Erkennungspipeline}
\subsection{Palm Detection}
\subsection{Handlandmarks}
\subsection{ROI}
\subsubsection{Lifecycle}

\section{Datenaugmentation und Modellquantisierung}
\subsection{Datenaugmentation}
\subsection{Modellquantisierung}

\section{Edge AI und Neuronale Beschleuniger (NPU)}
\subsection{Edge AI}
\subsection{Neuronale Beschleuniger (NPU)}
\subsection{ST Edge Core CLI}
